\section{Moment Amplification}

% Comparison with Moment Estimation
Instead of measuring the phases of the effective Grover operator, we now investigate the effect of uncontrolled Grover operations on the distributions.
We find an analogon in discrete increases of the canonical parameters when the measurement distribution of the start state is in an exponential family.


%We here apply \computationCircuits{} on anti-symmetric ancilla states (as for sign encoding)
%\begin{itemize}
%    \item A single \computationCircuit{} effectively prepares a reflection across the models of the computed formula
%    \item Building the reflections based on ancilla ground states (as for basis encoding), we would need an additional Pauli-Z and two \computationCircuits{}
%    \item Further ancilla qubits used in circuit decompositions need to be uncomputed (i.e. by applying the same partial \computationCircuit{} again) for further usage
%\end{itemize}

\subsection{Effect of Grover rotations}



%\subsection{Rotation along a formula}

Rotating $\repnum$ times across the initial state then gives the tensor product of the antisymmetric ancilla state with
\begin{align*}
    \qstateofat{\repnum}{\shortcatvariables}
    = \sinof{\left(\frac{1}{2}+\repnum\right)\rotanglesymbol} \qstateofat{\probof{0},\exformula}{\shortcatvariables} +
    \cosof{\left(\frac{1}{2}+\repnum\right)\rotanglesymbol} \qstateofat{\probof{0},\lnot\exformula}{\shortcatvariables} \, .
\end{align*}
This reflection can be performed by the operator on the distributed qubits
\begin{align*}
    U^t \left(
            \identityat{\shortcatvariables^{\insymbol},\shortcatvariables^{\outsymbol}}-2\onehotmapofat{0}{\shortcatvariables^{\insymbol}} \otimes \onehotmapofat{0}{\shortcatvariables^{\outsymbol}}
    \right) U \, .
\end{align*}


\subsubsection{Mean parameter interpretation}
% Mean parameter interpretation

% As exponential family parameters
The mean parameter of $\exformula$ with respect to the measurement distribution is
\begin{align*}
    \meanparam^{\repnum}
    = \contraction{\probofat{\repnum}{\shortcatvariables},\formulawith}
    = \sinof{\left(\frac{1}{2}+\repnum\right)\rotanglesymbol}^2
    = \sinof{\left(1+2\repnum\right)\sin^{-1}\left(\sqrt{\meanparam^{0}}\right)}^2
    %= \left(\frac{
    %    2\cdot \sinof{\left(\frac{1}{2}+\repnum\right)\rotanglesymbol} \cdot \sin^{-1}\left(\sqrt{\meanparam^{0}}\right)
    %}{\meanparam^{0}}\right)^2
\end{align*}
Here we used that $\rotanglesymbol = 2\cdot\sin^{-1}\left(\sqrt{\meanparam^{0}}\right)$.


\subsubsection{Canonical parameter interpretation}

%\begin{figure}[t]
%    \begin{center}
%        \input{partII/5-quantum-inference/tikz/moment-amplification/rotation-sketch}
%    \end{center}
%    \caption{Geometric visualization of the states $\qstateof{1}$ resulting from a Grover rotation on $\qstateof{0}$, in the subspace spanned by $\qstateof{\probof{0},\exformula}$ and $\qstateof{\probof{0},\lnot\exformula}$.
%    The effective reflection by the \computationCircuit{} results in the gray state $\qstateof{\badstatesymbol}-\qstateof{\goodstatesymbol}$.
%    Reflecting again on the initial state $\qstateof{0}$ results in the amplified state $\qstateof{1}$.
%    Each such Grover iteration is therefore a rotation by the angle $\rotanglesymbol$ in the.}
%\end{figure}



\begin{lemma}
    \label{lem:rotationCanParamIncrease}
    \red{Works also for more generic measurement distributions!}
    Let $\qstateof{0}$ be the q-sample of the elementary \ComputationActivationNetwork{} with canonical parameters $\hybridparam$, and let us amplify the formula $\formulaof{\selindex}$.
    Then $\qstateof{\repnum}$ is the q-sample of the elementary \ComputationActivationNetwork{} with canonical parameters
    \begin{align*}
        \hardlegset^{\repnum} =
        \begin{cases}
            \hardlegset \cup \{\selindex\}, \headindexof{\hardlegset}^{\selindex}
            = (\headindexof{\hardlegset},1) & \ifspace \sinof{(1+2\cdot\repnum)\sin^{-1}\left(\sqrt{\meanparamofat{0}{\indexedselvariable}}\right)} = 1 \\
            \hardlegset, \headindexof{\hardlegset} & \elsetext
        \end{cases}
    \end{align*}
    and
    \begin{align*}
        \canparamofat{\repnum}{\indexedselvariable} =
        \begin{cases}
            0 & \ifspace  \sinof{(1+2\cdot\repnum)\sin^{-1}\left(\sqrt{\meanparamofat{0}{\indexedselvariable}}\right)} = 1 \\
            \canparamat{\indexedselvariable} + \frac{
                \cos^2\left(\frac{\rotanglesymbol}{2}\right) \left(1-
                                                                 \cos^2\left(\left(\frac{1}{2}+\repnum\right) \rotanglesymbol \right)
                \right)
            }{
                \sin^2\left(\frac{\rotanglesymbol}{2}\right)
                \cdot \cos^2\left(\left(\frac{1}{2}+\repnum\right)\rotanglesymbol\right)
            } & \elsetext
        \end{cases} \, .
    \end{align*}
\end{lemma}
\begin{proof}
    Notice, that
    \begin{align*}
        \contractionof{\probofat{0}{\shortcatvariables},\bencodingofat{\formulaof{\selindex}}{\headvariableof{\selindex},\shortcatvariables}}{\headvariableof{\selindex}}
        = \coloredmatrixof{
            \left(\cosof{\frac{\rotanglesymbol}{2}}\right)^2 \\
            \left(\sinof{\frac{\rotanglesymbol}{2}}\right)^2
        }{\headvariableof{\selindex}}
        = \coloredmatrixof{
            1-\meanparamofat{0}{\indexedselvariable} \\
            \meanparamofat{0}{\indexedselvariable}
        }{\headvariableof{\selindex}}
    \end{align*}
    and
    \begin{align*}
        \contractionof{\probofat{\repnum}{\shortcatvariables},\bencodingofat{\formulaof{\selindex}}{\headvariableof{\selindex},\shortcatvariables}}{\headvariableof{\selindex}}
        =
        \coloredmatrixof{
            \left(\cosof{\left(\frac{1}{2}+\repnum\right)\rotanglesymbol}\right)^2 \\
            \left(\sinof{\left(\frac{1}{2}+\repnum\right)\rotanglesymbol}\right)^2
        }{\headvariableof{\selindex}}
        \coloredmatrixof{
            1-\meanparamofat{\repnum}{\indexedselvariable} \\
            \meanparamofat{\repnum}{\indexedselvariable}
        }{\headvariableof{\selindex}}
    \end{align*}
    If $\meanparamofat{\repnum}{\indexedselvariable}=1$ (that is $\sinof{\left(\frac{1}{2}+\repnum\right)\rotanglesymbol}=1$), the amplitude amplification amounts to adding $\formulaof{\selindex}$ as a hard constraint.
    This is equal to $\hardlegset^{\repnum}=\hardlegset\cup\{\selindex\}$ and $\headindexof{\hardlegset}^{\repnum}=(\headindexof{\hardlegset},1)$.

    If $\meanparamofat{\repnum}{\indexedselvariable}\neq 1$ we choose $\canparamofat{\Delta}{\indexedselvariable}\geq 0$ as in the claim and have after some algebra
    \begin{align*}
        \frac{
            \left(\cosof{\frac{\rotanglesymbol}{2}}\right)^2
        }{
            \left(\cosof{\frac{\rotanglesymbol}{2}}\right)^2 + \canparamofat{\Delta}{\indexedselvariable} \cdot \left(\sinof{\frac{\rotanglesymbol}{2}}\right)^2
        } = \left(\cosof{\left(\frac{1}{2}+\repnum\right)\rotanglesymbol}\right)^2
    \end{align*}
    and therefore
    \begin{align*}
        \contractionof{\probofat{\repnum}{\shortcatvariables},\bencodingofat{\formulaof{\selindex}}{\headvariableof{\selindex},\shortcatvariables}}{\headvariableof{\selindex}}
        = \normalizationof{
            \expof{\canparamofat{\Delta}{\indexedselvariable}\cdot\formulaofat{\selindex}{\canparam}}
            \probofat{0}{\shortcatvariables},\bencodingofat{\formulaof{\selindex}}{\headvariableof{\selindex},\shortcatvariables}
        }{\headvariableof{\selindex}} \, .
    \end{align*}
    Thus, the rotation of the q-sample corresponds with an increase of the canonical parameter $\canparamofat{0}{\indexedselvariable}$ by $\canparamofat{\Delta}{\indexedselvariable}$.
\end{proof}

\subsubsection{Maximum Entropy}

% Q-sample
When $\left(\frac{1}{2}+\repnum\right)\rotanglesymbol\leq \frac{\pi}{2}$ then all amplitudes remain positive, and $\qstateofat{\repnum}{\shortcatvariables}$ is a q-sample.
What is more, we can show that the property of maximum entropy with respect to a statistic containing the amplified formula.

\begin{theorem}
    \label{the:rotationKeepsMaxEntropy}
    Let $\qstateof{0}$ be the q-sample of a maximum entropy distribution with respect to the uniform base measure and the statistic $\formulaset$ containing the amplified formula, and let $\repnum\in\nn$ be a rotation number such that $\left(\frac{1}{2}+\repnum\right)\rotanglesymbol\leq\frac{\pi}{2}$.
    Then also $\qstateof{\repnum}$ is the q-sample of a maximum entropy distribution.
\end{theorem}
%We show \theref{the:rotationKeepsMaxEntropy} later based on a more technical description of the rotated state, which we show as the next lemma.
\begin{proof}%[Proof of \theref{the:rotationKeepsMaxEntropy}]
    By \lemref{lem:rotationCanParamIncrease} the amplitude amplified state is an elementary \ComputationActivationNetwork{} and thus a maximum entropy distribution.
        %   Since the measurement distribution of any quantum state in the plane spanned by $\qstateof{\goodstatesymbol}$ and $\qstateof{\badstatesymbol}$ is an elementary \ComputationActivationNetwork{}.
        %   When the assumption on $\repnum$ is met, we further have positive real amplitudes and therefore a q-sample.
        %  Since each amplitude amplification can be understood as an change of mean parameters.
        %  To be more precise, we characterize the change of canonical parameters in the lemma below.
\end{proof}

%
Notice, that we can only increase the canonical parameter by amplitude amplification.
Decreasing could be done by amplitude amplification on $\lnot\formulaof{\selindex}$ instead.



\subsection{Moment Matching} %based preparation of \CompActNets{}

We now investigate which \ComputationActivationNetworks{} can be prepared directly with amplitude amplification.
While we already have studied sparse representations of the computation network parts, we here apply amplitude amplification by Grover rotations as an activation mechanism.

We orient on the iterative proportional fitting algorithm, which is used to estimate the canonical parameters in exponential families  \cite{wainwright_graphical_2008} and more generally of \ComputationActivationNetworks{} \cite{goessmann_tensor-network_2025}.

The algorithm loops over the moment estimation and moment amplification subroutines, which have been introduced in the sections above.
The idea is to first compute the current moments of a distribution and second amplify one.



%This change of mean parameters can be understood by increasing the canonical parameter to $\exformula$ by
%\begin{align*}
%    \canparamof{\repnum} - \canparamof{0}
%    = \frac{
%        \cos^2\left(\frac{\rotanglesymbol}{2}\right) \left(1-
%                                                         \cos^2\left(\left(\frac{1}{2}+\repnum\right) \rotanglesymbol \right)
%        \right)
%    }{
%        \sin^2\left(\frac{\rotanglesymbol}{2}\right)
%        \cdot \cos^2\left(\left(\frac{1}{2}+\repnum\right)\rotanglesymbol\right)
%    }
%\end{align*}
%we here allow for $\canparam=\infty$, which corresponds to the case of a hard activation to $\exformula$.

We understand this rotation as a moment matching operation to the formula $\exformula$, see \algoref{alg:QCMM}.

\begin{algorithm}
    \caption{Quantum Circuit Moment Matcher}\label{alg:QCMM}
    \begin{algorithmic}
        \Require Formulas $\formulaofat{\seldim}{\shortcatvariables}$ for $\selindexin$, vector $\datameanat{\selvariable}$ of mean parameters
        \Ensure Circuit $\unitarysymbol$ preparing a state (when applied on the ground state) which measurement distribution matched $\datameanat{\selvariable}$\\
        \hrule
        \State $U\algdefsymbol (\paulixsymbol\circ\hgate)\left[\avariableof{\insymbol},\avariableof{\outsymbol}\right] \otimes
        \left(\bigotimes_{\catenumeratorin}\hgateat{\catvariableof{\catenumerator,\insymbol},\catvariableof{\catenumerator,\outsymbol}}\right)
        $
        \While{Convergence criterion not met}
%        \For{$\selindexin$}
            \State Estimate $\meanparamat{\selvariable}$,
            \begin{itemize}
                \item either by particle bases inference
                \item or by quantum counting (measuring the eigenvalue of the Grover operator), see \secref{sec:amplitudeEstimation}
            \end{itemize}
            \State Select an $\selindexin$, which moment has to be updated (e.g. by comparison of $\meanparamat{\selvariable}$ with $\datameanat{\selvariable}$)
            \State Choose whether to amplify $\formulaof{\selindex}$ or $\lnot\formulaof{\selindex}$ based on whether the estimated mean is smaller than the target mean
            \State Compute optimal $\repnum$ to match $\datameanat{\selvariable}$
            \State Extend Circuit $\unitarysymbol$ by $\repnum$ alternations of
            \begin{itemize}
                \item $\qcbencodingof{\formulaof{\selindex}}$ (effectively an reflection across the subspace spanned by one-hot encoded models)
                \item reflections across the current state (by $U(I-2e_0e_0)U$)
            \end{itemize}
        \EndWhile
    \end{algorithmic}
\end{algorithm}


% Outook
We can use \algoref{alg:QCMM} as a preparation algorithm of a q-sample before doing ancilla augmentation as in the main section.
Note that is important to keep track of the canonical parameters in the preparation.
When activating the statistic qubits, only the difference of the already prepared canonical parameters and the target canonical parameters has to be used (importance sampling interpretation).

